\cchapter{Scripting and animations}

The computations defined in chapter \ref{Chapter:computations} can
also be performed in batch mode. This is useful when a large number
of files need to be manipulated at once. Two other powerful features
of \TexMol include
\begin{itemize}
\item Scripting.
\item Animation.
\end{itemize}
In this chapter, we will go through simple examples for each of the
above three features.

\csection{Batch-mode calculations}

Electron density, hydrophobicity and curvatures can be estimated
from the command line. The commands are as follows.

\subsection{Electron density and hydrophobicity}

The command line for creating an electron density or hydrophobicity
map is:

\TexMol \textbf{-blur} \param{input molecule file name}
\param{output volume file name} \param{dim1} \param{dim2}
\param{dim3}
\param{density type} \param{colored volume} \param{gaussian
blobbiness}
\param{color} \param{colormap file name} \param{gap}

The different parameters are:

\param{input molecule file name} The full path and name of the input molecule
file.

\param{output volume file name} The full path and name of the output
rawiv or rawv file.

\param{dim1, dim2, dim3} Dimensions of the output volume.

\param{density type} Enter \mm{0} for electron density and \mm{1} for
hydrophobicity.

\param{colored volume} \emph{true} if you want a rawv volume.

\param{gaussian blobbiness} We recommend a value of -2.3 to conform
with previous research and bigger values, say -0.1, to get smoother
shapes representing the volume at much lower resolutions.

\param{color} If the colored volume parameter was true, then we can
either specify either color by structure or provide a color map
file. If we do not provide a color map file, then this parameter is
relevant. \mm{0} stands for atom coloring, \mm{1} for residue,
\mm{2} for secondary structure and \mm{3} for chain coloring.

\param{colormap file name} The molecule can be colored according to
some user defined colors, to perhaps highlight some structure which
cannot be done in the more general method previous described. The
file format for the \emph{colormap file name} is given in the file
format description page from CCV's website.

\param{gap} This should be used sparingly, to provide a gap of empty
space around the requested dimensions. It is best left as \mm{0}.

\subsection{Curvatures}

There are two commands for generating curvatures. One is from a
volume file with a given isovalue, where \TexMol extracts the mesh
where we need to obtain the curvature, and second, where the user
inputs the mesh directly. The two commands in order are:

\TexMol \textbf{-setcurvature} \param{1} \param{input molecule file
name} \param{output volume file name} \param{output mean curvature
file name} \param{output gaussian curvature file name} \param{dim1}
\param{dim2} \param{dim3} \param{gaussian
blobbiness} \param{isovalue} \param{number of grid subdivisions}
\param{maximum function error}\\

\TexMol \textbf{-setcurvature} \param{0} \param{input molecule file
name} \param{input surface file name} \param{output mean curvature
file name} \param{output gaussian curvature file name} \param{dim1}
\param{dim2} \param{dim3} \param{gaussian
blobbiness} \param{number of grid subdivisions}
\param{maximum function error}\\

The various parameters are described below.

\param{0} and \param{1} differentiate between the two calls.

\param{input molecule file name} The full path and name of the input molecule
file.

\param{output volume file name} The full path and name of the output
rawiv or rawv file.

\param{output mean curvature file name} The full path and name of
the output mean curvature file.

\param{output gaussian curvature file name} The full path and name of
the output gaussian curvature file.

\param{dim1, dim2, dim3} Dimensions of the output volume.

\param{gaussian blobbiness} We recommend a value of -2.3 to conform
with previous research and bigger values, say -0.1, to get smoother
shapes representing the volume at much lower resolutions.

\param{number of grid subdivisions} The higher this value, the more
the precomputation and memory requirement. There is a cut off value
around 10 to 20 where you get optimum speed for a given machine.

\param{maximum function error} The maximum allowable function error
while calculating the curvatures.

\warning{While creating curvatures, if we overestimate the value of
grid spacing, we could end up with very high memory usages and may
have to kill the program.}


\csection{Scripting}

\TexMol currently offers a limited but powerful set of functions as
part of its scripting module. This module is available once \TexMol
is run in user interface mode. The parser command window can be
opened from Utilities - script from the main menu bar. Here users
can enter commands, select one or more and execute them. Some of the
commands available are listed below.

\textcolor{BrickRed}{bool} blur(\textcolor{BrickRed}{int} argc,
\textcolor{BrickRed}{char*} argv[]);

\textcolor{BrickRed}{bool} setCurvature(\textcolor{BrickRed}{int}
argc, \textcolor{BrickRed}{char*} argv[]);

\textcolor{BrickRed}{bool}
outGridPositions(\textcolor{BrickRed}{int} argc,
\textcolor{BrickRed}{char*} argv[]);

\textcolor{BrickRed}{bool} classifyPoints(\textcolor{BrickRed}{int}
argc, \textcolor{BrickRed}{char*} argv[]);

\textcolor{BrickRed}{bool} growOut(\textcolor{BrickRed}{int} argc,
\textcolor{BrickRed}{char*} argv[]);

\textcolor{BrickRed}{bool} getSurface(\textcolor{BrickRed}{int}
argc, \textcolor{BrickRed}{char*} argv[]);

\textcolor{BrickRed}{bool} evolve(\textcolor{BrickRed}{int} argc,
\textcolor{BrickRed}{char*} argv[]);

\textcolor{BrickRed}{bool} writePDB(\textcolor{BrickRed}{int} argc,
\textcolor{BrickRed}{char*} argv[]);

\textcolor{BrickRed}{bool} writeGOA(\textcolor{BrickRed}{int} argc,
\textcolor{BrickRed}{char*} argv[]);

\textcolor{BrickRed}{bool} depthColor(\textcolor{BrickRed}{int}
argc, \textcolor{BrickRed}{char*} argv[]);

\textcolor{BrickRed}{bool}
printPDBInformation(\textcolor{BrickRed}{int} argc,
\textcolor{BrickRed}{char*} argv[]);

\textcolor{BrickRed}{bool}
getMaxDistanceFromPoint(\textcolor{BrickRed}{int} argc,
\textcolor{BrickRed}{char*} argv[]);

\textcolor{BrickRed}{bool} addNewDataSet(\textcolor{BrickRed}{int}
argc, \textcolor{BrickRed}{char*} argv[],
\textcolor{BrickRed}{MoleculeVizMainWindow} *mWindow );

\textcolor{BrickRed}{bool} splitView(\textcolor{BrickRed}{int} argc,
\textcolor{BrickRed}{char*} argv[],
\textcolor{BrickRed}{MoleculeVizMainWindow} *mWindow );

\textcolor{BrickRed}{bool} deleteData(\textcolor{BrickRed}{int}
argc, \textcolor{BrickRed}{char*} argv[],
\textcolor{BrickRed}{MoleculeVizMainWindow} *mWindow );


\textcolor{BrickRed}{bool} deletePrevData(\textcolor{BrickRed}{int}
argc, \textcolor{BrickRed}{char*} argv[],
\textcolor{BrickRed}{MoleculeVizMainWindow} *mWindow );


\textcolor{BrickRed}{bool} deleteAllData(\textcolor{BrickRed}{int}
argc, \textcolor{BrickRed}{char*} argv[],
\textcolor{BrickRed}{MoleculeVizMainWindow} *mWindow );


\textcolor{BrickRed}{bool} setVisible(\textcolor{BrickRed}{int}
argc, \textcolor{BrickRed}{char*} argv[],
\textcolor{BrickRed}{MoleculeVizMainWindow} *mWindow );


\textcolor{BrickRed}{bool} setVisiblePrev(\textcolor{BrickRed}{int}
argc, \textcolor{BrickRed}{char*} argv[],
\textcolor{BrickRed}{MoleculeVizMainWindow} *mWindow );


\textcolor{BrickRed}{bool} saveImage(\textcolor{BrickRed}{int} argc,
\textcolor{BrickRed}{char*} argv[],
\textcolor{BrickRed}{MoleculeVizMainWindow} *mWindow );


\textcolor{BrickRed}{bool} setGridVisible(\textcolor{BrickRed}{int}
argc, \textcolor{BrickRed}{char*} argv[],
\textcolor{BrickRed}{MoleculeVizMainWindow} *mWindow );



\csection{Animations}

This module is mainly useful for making movies. Animations can be
saved as a sequence of images in \TexMol. We allow the user to load
a data set, render it and transform the view point, and save
trajectories and play it back. There are actually two steps in the
animation process.

\begin{enumerate}
\item Create and save a trajectory.
\item Save the animation based on a trajectory.
\end{enumerate}

\subsection{Creating a trajectory}
\label{subsection:CreateTrajectory}

Data sets can be quite large, making them hard to interact with.
Hence, it is sometimes useful to save a trajectory using a low
resolution data set and later use it to render a large data set in
high resolution. In order to create the trajectory it is useful to
have a low resolution data set which the user can comfortably
rotate, zoom and translate interactively. To save a trajectory,
follow the procedure outlined in table \ref{Table:Trajectory}

\steps{Creating and saving a trajectory}{Table:Trajectory}{\item
Load and display the data set or set of data sets you want to create
a movie of.
\item Use the mouse to move them into the correct starting viewing
position. \item Click Animation, Start recording from the menu bar.
\item Enter and save the file name of the trajectory file. It does
not require any specific file extension. \item On pressing OK, the
recording is on. Any mouse movement in the rendering area is saved
in to the file, along  with the time. \item Once you are done with
transforming the data, click Animation, Stop recording in the menu
bar. \item You can see if the trajectory was satisfactory by playing
it back by clicking on Animation, Playback animation from the menu
bar.}

\subsection{Saving animations}

Before you can save a sequence of images, you need to create a
trajectory as described in section
\ref{subsection:CreateTrajectory}.


\steps{Saving an animation}{Table:SavingAnimations}{ \item Load the
high resolution data sets to save.
\item keep the screen resolution and window size as high as
required.
\item Open the trajectory file by selecting Animation, Record
animation from the menu bar. \item \TexMol will automatically begin
to replay the trajectory and save images.}

\warning{Disable the screensaver as saving images in high resolution
can take a long time. We also interpolate between mouse movements
using linear interpolation to get better smoothness, causing the
movie to be longer than expected.}
